\chapter{What have we learned?}
\label{Chpt:Summary}

In this chapter we summarize the main results and lessons learned from the Dy-K experiment. In the next chapter, we will instead provide an outlook on future studies achievable with our experiment, and generalize to the study of mass-imbalanced Fermi gases and to unconventional quantum mixtures.

The first important lesson drawn from the experiment, which is probably not (yet) fully conveyed in the publications from our group, is the difficulty of producing the doubly degenerate mixture. While the evaporative cooling of dysprosium and the simultaneous sympathetic cooling of potassium are reliable and achieve good degeneracy in both species, satisfactory cooling is achieved only at very low magnetic fields, below \SI{500}{mG}. Above this magnetic field, the interaction spectrum of Dy is observed to be strongly temperature-dependent, and severe three-body losses prevent efficient evaporative cooling.  In a first study of the interaction spectrum of $^{161}$Dy, see Ref.~\cite{Soave2022lff}, numerous intraspecies resonances were identified for a small magnetic field range, but a complete characterization of these resonances is still lacking. A similar behavior of the intraspecies resonances caused by thermal effects  was identified in a sample of $^{166}$Er \cite{Krstajic2025cot}; however, the density of resonances for the case of bosonic erbium is much lower than for the case of our fermionic dysprosium. These limitations, for a long time not fully explored in order to characterize interactions in the mixture, are currently being studied in our experiment, with the aim of better understanding and mitigating these effects.


A direct consequences of the lack of evaporative cooling for high magnetic fields is the necessity of quickly switching magnetic field, from the preparation region to the vicinity of a Feshbach resonance. The necessity of ramping the magnetic field over several \unit{G}, up to hundreds,  

One of the partial work around developed in this thesis to this issue concern the identification of several resonances in the low magnetic field region, closer to the preparation field. 


Over reliance on the magnetic field gradient: great useful tool, that revealed to be at the same time a limitation, especially for the study of Feshbach molecules.

Metallic chamber waay too big and adding too many eddy currents, solution to be combined to the imaging one, have a glass cell in which to do everything.

Poor imaging, fine for what was the goal up to now, but since early days one of the main limiting factors. a long transport towards a different chamber could be challenging for the wavelengths available a full glass cell solution, like the one implemented in the NewlanD experiment of Dr. Emil Kirilov in our group, has the potential to be a more suitable platform.


Train of toughs of this chapter
no preparation at different fields -> a lot of ramps (how we try to fix some of the aspects)-> magnetic field gradients has resource but also as a limitation -> also discovery of light induced losses, poor coating to plan around -> low resolution imaging -> second generation exp needed to fix many of these short comings now that we have a much more complete picture of the system and what would be necessary to push its limits. the current studies in the experiments are dedicated towards understanding the final limitation in order to have a complete set of what is necessary for the future step and reaching superfluidity in our (or similar systems)

